A real quadratic space with nondegenerate part of signature $(p,q)$ and radical of dimension $r$ admits a nondegenerate completion; every completion has dimension at least $p + q + 2r$; the minimal ones have signature exactly $(p+r, q+r)$; and any two minimal completions are isometric by an isometry fixing the original space pointwise, so the minimal ambient Clifford algebra $\Cl(p+r, q+r)$ is unique up to an isomorphism fixing the embedded $\Cl(p,q;r)$. In the model completion the radical is paired with its dual, $Q(\rho, \varphi) = 1$, $Q(\varphi,\varphi) = 0$, $\varphi$ orthogonal to the nondegenerate part. For a type-N obstruction in a nondegenerate algebra the increment $(+1,+1)$ is forced outright; for a type-C obstruction the completed algebra is unique but the sign of the adjoined generator's square is a geometric datum the algebra forgets. \emph{Provenance:} \cite[Thm.~2.3 and its model]{ButlerT2} (\claimref{conformal.minimal-nondegenerate-completion}); \cite[Thm.~4.13, Thm.~5.7 and Rem.~5.8]{ButlerT1} (\claimref{transport.minimality}, \claimref{transport.one-generator-minimal-minimal-completions-classified}).
